
\begin{abstract}
Developers often evolve an existing software system by making
internal changes, called \emph{migration}. Moving to a new framework,
changing implementation to improve efficiency, and upgrading a
dependency to its latest version are examples of migrations.

Migration is a common and typically continuous maintenance
task undertaken either manually or through tooling. Certain 
migrations are labor intensive and costly, developers do not find the
required work rewarding, and they may take years to complete.
Hence, automation is preferred for such migrations.

In this paper, we discuss a large-scale, costly and traditionally
manual migration project at Google, propose a novel automated
algorithm that uses change location discovery and a Large Language
Model (LLM) to aid developers conduct the migration, report the
results of a large case study, and discuss lessons learned.

Our case study on 39 distinct migrations undertaken by three developers 
over twelve months shows that a total of 595 code changes
with 93,574 edits have been submitted, where 74.45\% of the code
changes and 69.46\% of the edits were generated by the LLM. The 
developers reported high satisfaction with the automated tooling, and
estimated a 50\% reduction on the total time spent on the migration
compared to earlier manual migrations.

Our results suggest that our automated, LLM-assisted workflow
can serve as a model for similar initiatives.
\end{abstract}

\keywords{Software, Migration, Refactoring, Transformation, Productivity, LLM}
